\documentclass[10pt]{article}
\usepackage[letterpaper,margin=0.6in]{geometry}
\usepackage{amsmath,amssymb,booktabs,tabularx,array}
\usepackage[hidelinks]{hyperref}
\pagestyle{empty}
\setlength{\parindent}{0pt}\setlength{\parskip}{3pt}
\newcolumntype{L}{>{\raggedright\arraybackslash}X}
\newcolumntype{P}[1]{>{\raggedright\arraybackslash}p{#1}}
\begin{document}
{\Large\bfseries One-sheet: the Feynman path sum at closure order $\alpha=2$}\par
\textit{An Axiomatic Derivation of the Feynman Path Sum at the Light-Cone Kernel Width and Closure Order $\alpha=2$ of the Causal Accumulation Law.} S.~Leizerman, LIOR Research Initiative, November 13, 2026. Script: \texttt{cal\_feynman\_alpha2.py} (numpy only, about 25\,s). Units: $\hbar=c=1$, $\tau_c=1/\sqrt2$, so $\mu=1$.

\textbf{Inputs.}
\begin{itemize}\setlength{\itemsep}{0pt}
\item The CAL at $\alpha=2$: $T=2J+K_2J$. The transport chain is $1$, so there is no coupling constant.
\item A light-cone kernel: at lag $\tau$ the profile is the light front of radius $c\tau$.
\item The balance reading: a free carrier balances its light-cone past ($T=0$), and a source is the imbalance. The edge's update rule is not used here.
\item The phase belongs to the source. The double accumulation gives the half turn, and no factor goes on the memory.
\item For the readout, one sign of frequency, inherited from the source's past.
\end{itemize}
The kernel's lag scale $\tau_c$ is the only scale; only ratios such as $\mu\tau_c=1/\sqrt2$ are physical.

\medskip
\begin{tabularx}{\textwidth}{@{}P{3.3cm}LP{4.4cm}@{}}
\toprule
\textbf{Result} & \textbf{Statement} & \textbf{Check}\\
\midrule
Lorentz selects $\alpha=2$ & The light-cone symbol depends on $s^2+c^2k^2$ alone only at $\alpha=2$ (and $0$). The defect is $\alpha(2-\alpha)\varphi^2/6$. & defect $0$ at $\alpha=0,2$; $0.10$ at $1.9$, $2.1$\\
Kernel $=$ Green's function & Kirchhoff's shell in $3+1$; one order-$1$ accumulation per null leg in $1+1$ ($2=1+1$). & closed form matches direct integration to 6 digits\\
Klein--Gordon & The balanced carrier obeys $(\square+\mu^2)J=0$ with $\mu\tau_c=1/\sqrt2$. The mass is the kernel's own rate. & roots $=\sqrt{k^2+\mu^2}$ to 10 digits\\
Conservative at $2$ & Both signs decay at $\tfrac\pi4(2-\alpha)\mu^2/\omega$ below $2$ and run away at that rate above. Nothing near $2$ picks a sign. & ${\rm Im}\,\omega$: $-5.5\times10^{-3}$, $0$, $+5.6\times10^{-3}$ at $1.99$, $2$, $2.01$\\
Path sum & The Neumann series is a sum over piecewise-null histories with $-\mu^2$ per corner. It converges absolutely and resums to $G_\mu$. & $1+1$ residual $3.35\times10^{-5}\to8.38\times10^{-6}$\\
Checkerboard & The $n$-th term is the $(2n+1)$-corner amplitude over $2ic\mu$, since $(-\mu^2)^n=(i\mu)^{2n}$. & lattice error $\propto\varepsilon$\\
Corner tail & In $3+1$, histories without corners sit on the cone; the corners alone build the $J_1$ tail inside it. & $\le2.5\times10^{-6}$\\
Unitarity, readout & Each sign of frequency is unitary. On one sign, $\int JJ^\dagger$ is conserved, so the readout is Born-normalized in every frame. & spread $2.5\times10^{-16}$ vs $1.6\times10^{-1}$\\
Feynman's kernel & The positive-frequency path sum, rest phase removed, gives $K_F$ as $c\to\infty$ with error $O(1/c^2)$. & error$\cdot c^2\to3.667$ (pred.\ $3.672$)\\
\bottomrule
\end{tabularx}

\medskip
\textbf{Where the $i$ comes from.} The $\alpha=2$ dynamics is real. The $i$ of $K_F$ enters only with the sign of frequency inherited from the source. At the edge it is the rotation produced by negative feedback.

\textbf{Not claimed / open.}
\begin{itemize}\setlength{\itemsep}{0pt}
\item Interaction between carriers (the return leg).
\item Spin (a Dirac form in $3+1$).
\item Mass ratios between carriers.
\item The relation to the edge's rest rotation; the equivalence condition is a forward pointer only.
\item A frame-independent position readout (Newton--Wigner).
\item Prior art: the $1+1$ checkerboard (Feynman--Hibbs, Jacobson--Schulman, GJKS) and causal-set chain sums (Johnston). What is new is the source: the history sum is the CAL's memory series.
\end{itemize}
\end{document}
